\documentclass[12pt,letterpaper]{article}

%------------------------------------------------
% PAQUETES
%------------------------------------------------
\usepackage[spanish]{babel}
\usepackage[utf8]{inputenc}
\usepackage[T1]{fontenc}

\usepackage{graphicx}
\usepackage{geometry}
\usepackage{fancyhdr}
\usepackage{amsmath,amssymb}
\usepackage{xcolor}
\usepackage{enumitem}

%------------------------------------------------
% CONFIGURACIÓN DE PÁGINA
%------------------------------------------------
\geometry{margin=2.5cm}
\setlength{\headheight}{52pt}

%------------------------------------------------
% ENCABEZADO
%------------------------------------------------
\pagestyle{fancy}
\fancyhf{}

\fancyhead[L]{
    \includegraphics[height=1.5cm]{ur_image.png}
}

\fancyhead[R]{
\begin{tabular}{r}
\textbf{Monitoría Álgebra Lineal} \\
Gabriel Fernando Salcedo Zamora
\end{tabular}
}

\renewcommand{\headrulewidth}{0.4pt}

%------------------------------------------------
% COMANDO PARA RESPUESTAS
%------------------------------------------------
\newcommand{\respuesta}[1]{
\[
\textcolor{blue}{#1}
\]
}

%------------------------------------------------
% DOCUMENTO
%------------------------------------------------
\begin{document}

\begin{center}
    {\LARGE \textbf{Gram--Schmidt}}\\[0.3cm]
    \rule{\textwidth}{0.4pt}
\end{center}

%================================================
% EJERCICIOS
%================================================

\section*{Ejercicios}

\begin{enumerate}[label=\textbf{\arabic*.}]

%------------------------------------------------
\item Sean los vectores

\[
v_1=(1,0),
\qquad
v_2=(1,1)
\]

\begin{enumerate}[label=\alph*)]

\item Tome
\[
u_1=v_1
\]

\item Calcule la proyección de \(v_2\) sobre \(u_1\).

\item Halle
\[
u_2=v_2-\operatorname{proj}_{u_1}(v_2)
\]

\item Verifique que \(u_1\) y \(u_2\) sean ortogonales.

\item Normalice los vectores obtenidos.

\end{enumerate}

\vspace{0.5cm}

%------------------------------------------------
\item Aplique el proceso de Gram--Schmidt al conjunto

\[
v_1=(1,1),
\qquad
v_2=(1,-1)
\]

\begin{enumerate}[label=\alph*)]

\item Verifique si los vectores ya son ortogonales.

\item Encuentre una base ortonormal.

\item Compruebe que cada vector tiene norma igual a 1.

\end{enumerate}

\vspace{0.5cm}

%------------------------------------------------
\item Considere los vectores

\[
v_1=(1,0,0),
\qquad
v_2=(1,1,0)
\]

\begin{enumerate}[label=\alph*)]

\item Calcule \(u_1\).

\item Halle la proyección de \(v_2\) sobre \(u_1\).

\item Determine \(u_2\).

\item Verifique que
\[
u_1\cdot u_2=0
\]

\item Obtenga una base ortonormal.

\end{enumerate}

\vspace{0.5cm}

%------------------------------------------------
\item Sean

\[
v_1=(1,1,0),
\qquad
v_2=(1,0,1),
\qquad
v_3=(0,1,1)
\]

\begin{enumerate}[label=\alph*)]

\item Aplique el proceso de Gram--Schmidt.

\item Encuentre los vectores ortogonales:
\[
u_1,\quad u_2,\quad u_3
\]

\item Verifique la ortogonalidad entre ellos.

\item Normalice únicamente el vector \(u_1\).

\end{enumerate}

\vspace{0.5cm}

%------------------------------------------------
\item Considere los vectores

\[
v_1=(1,1,1),
\qquad
v_2=(2,2,2)
\]

\begin{enumerate}[label=\alph*)]

\item Determine si los vectores son linealmente independientes.

\item Inicie el proceso de Gram--Schmidt.

\item Explique qué ocurre al calcular el segundo vector ortogonal.

\item Interprete geométricamente el resultado.

\end{enumerate}

\end{enumerate}

%================================================
% SOLUCIONES
%================================================

\newpage

\section*{Soluciones}

\begin{enumerate}[label=\textbf{\arabic*.}]

%================================================
% EJERCICIO 1
%================================================

\item Sean

\[
v_1=(1,0),
\qquad
v_2=(1,1)
\]

\begin{enumerate}[label=\alph*)]

%------------------------------------------------
\item Primer vector ortogonal

\respuesta{
u_1=v_1=(1,0)
}

%------------------------------------------------
\item Proyección de \(v_2\) sobre \(u_1\)

\respuesta{
\operatorname{proj}_{u_1}(v_2)
=
\frac{v_2\cdot u_1}{u_1\cdot u_1}u_1
}

\respuesta{
=
\frac{(1)(1)+(1)(0)}{(1)^2+(0)^2}(1,0)
}

\respuesta{
=(1)(1,0)=(1,0)
}

%------------------------------------------------
\item Cálculo de \(u_2\)

\respuesta{
u_2=v_2-\operatorname{proj}_{u_1}(v_2)
}

\respuesta{
=(1,1)-(1,0)=(0,1)
}

%------------------------------------------------
\item Verificación de ortogonalidad

\respuesta{
u_1\cdot u_2=(1)(0)+(0)(1)=0
}

Por lo tanto, los vectores son ortogonales.

%------------------------------------------------
\item Normalización

\respuesta{
\|u_1\|=\sqrt{1^2+0^2}=1
}

\respuesta{
\|u_2\|=\sqrt{0^2+1^2}=1
}

La base ortonormal es:

\respuesta{
\left\{
(1,0),
(0,1)
\right\}
}

\end{enumerate}

\vspace{0.5cm}

%================================================
% EJERCICIO 2
%================================================

\item Sean

\[
v_1=(1,1),
\qquad
v_2=(1,-1)
\]

\begin{enumerate}[label=\alph*)]

%------------------------------------------------
\item Verificación de ortogonalidad

\respuesta{
v_1\cdot v_2=(1)(1)+(1)(-1)=1-1=0
}

Los vectores ya son ortogonales.

%------------------------------------------------
\item Base ortonormal

\respuesta{
\|v_1\|=\sqrt{1^2+1^2}=\sqrt2
}

\respuesta{
\|v_2\|=\sqrt{1^2+(-1)^2}=\sqrt2
}

\respuesta{
e_1=
\left(
\frac1{\sqrt2},
\frac1{\sqrt2}
\right)
}

\respuesta{
e_2=
\left(
\frac1{\sqrt2},
\frac{-1}{\sqrt2}
\right)
}

%------------------------------------------------
\item Verificación de norma

\respuesta{
\|e_1\|=\|e_2\|=1
}

\end{enumerate}

\vspace{0.5cm}

%================================================
% EJERCICIO 3
%================================================

\item Sean

\[
v_1=(1,0,0),
\qquad
v_2=(1,1,0)
\]

\begin{enumerate}[label=\alph*)]

%------------------------------------------------
\item Primer vector ortogonal

\respuesta{
u_1=v_1=(1,0,0)
}

%------------------------------------------------
\item Proyección

\respuesta{
\operatorname{proj}_{u_1}(v_2)
=
\frac{v_2\cdot u_1}{u_1\cdot u_1}u_1
}

\respuesta{
=
\frac{1}{1}(1,0,0)
=(1,0,0)
}

%------------------------------------------------
\item Cálculo de \(u_2\)

\respuesta{
u_2=(1,1,0)-(1,0,0)
}

\respuesta{
u_2=(0,1,0)
}

%------------------------------------------------
\item Verificación

\respuesta{
u_1\cdot u_2=(1)(0)+(0)(1)+(0)(0)=0
}

%------------------------------------------------
\item Base ortonormal

Como ambos vectores tienen norma 1:

\respuesta{
\left\{
(1,0,0),
(0,1,0)
\right\}
}

\end{enumerate}

\vspace{0.5cm}

%================================================
% EJERCICIO 4
%================================================

\item Sean

\[
v_1=(1,1,0),
\qquad
v_2=(1,0,1),
\qquad
v_3=(0,1,1)
\]

\begin{enumerate}[label=\alph*)]

%------------------------------------------------
\item Aplicación de Gram--Schmidt

\respuesta{
u_1=v_1=(1,1,0)
}

\respuesta{
u_2=v_2-\operatorname{proj}_{u_1}(v_2)
}

\respuesta{
v_2\cdot u_1=1
}

\respuesta{
u_1\cdot u_1=2
}

\respuesta{
\operatorname{proj}_{u_1}(v_2)
=
\frac12(1,1,0)
=
\left(
\frac12,
\frac12,
0
\right)
}

\respuesta{
u_2=
\left(
1,0,1
\right)
-
\left(
\frac12,
\frac12,
0
\right)
}

\respuesta{
u_2=
\left(
\frac12,
-\frac12,
1
\right)
}

%------------------------------------------------
\item Cálculo de \(u_3\)

\respuesta{
u_3=
v_3
-
\operatorname{proj}_{u_1}(v_3)
-
\operatorname{proj}_{u_2}(v_3)
}

\respuesta{
\operatorname{proj}_{u_1}(v_3)
=
\left(
\frac12,
\frac12,
0
\right)
}

\respuesta{
\operatorname{proj}_{u_2}(v_3)
=
\left(
-\frac13,
\frac13,
-\frac23
\right)
}

\respuesta{
u_3=
\left(
-\frac16,
\frac16,
\frac53
\right)
}

%------------------------------------------------
\item Verificación de ortogonalidad

\respuesta{
u_1\cdot u_2=0
}

\respuesta{
u_1\cdot u_3=0
}

\respuesta{
u_2\cdot u_3=0
}

%------------------------------------------------
\item Normalización de \(u_1\)

\respuesta{
\|u_1\|=\sqrt2
}

\respuesta{
e_1=
\left(
\frac1{\sqrt2},
\frac1{\sqrt2},
0
\right)
}

\end{enumerate}

\vspace{0.5cm}

%================================================
% EJERCICIO 5
%================================================

\item Sean

\[
v_1=(1,1,1),
\qquad
v_2=(2,2,2)
\]

\begin{enumerate}[label=\alph*)]

%------------------------------------------------
\item Dependencia lineal

\respuesta{
v_2=2v_1
}

Por lo tanto, los vectores son linealmente dependientes.

%------------------------------------------------
\item Inicio de Gram--Schmidt

\respuesta{
u_1=v_1=(1,1,1)
}

\respuesta{
u_2=v_2-\operatorname{proj}_{u_1}(v_2)
}

%------------------------------------------------
\item Resultado del segundo vector

\respuesta{
\operatorname{proj}_{u_1}(v_2)=v_2
}

Entonces:

\respuesta{
u_2=(2,2,2)-(2,2,2)=(0,0,0)
}

%------------------------------------------------
\item Interpretación geométrica

Geométricamente, esto significa que \(v_2\) no aporta una nueva dirección, ya que apunta exactamente en la misma dirección de \(v_1\).

\end{enumerate}

\end{enumerate}

\end{document}
